% =============================================================================
% warn-graphics-bleed -- a figure that bleeds past the text block.
%
% A full-width figure drawn 7.4" wide (via \rule) sticks ~0.4" past the 7"
% block on one page. Graphics poking out are an honest layout slip, so a
% bleeding figure is a warning, not an error.
%
% Expected checker outcome: PASS with warning [text-block].
% =============================================================================
\documentclass[letterpaper,twocolumn,10pt]{article}
\usepackage{usenix}
\usepackage{lipsum}

\date{}
\title{\Large \bf A Submission with a Bleeding Figure}
\author{{\rm Alice Anonymous}\\One University \and {\rm Bob Blinded}\\Two Institute}

\begin{document}
\maketitle

\begin{abstract}
\lipsum[1][1-5]
\end{abstract}

\section{Introduction}
\lipsum[2-3]

\begin{figure*}[t]
  \centering
  \makebox[0pt]{\rule{7.4in}{1.2in}}
  \caption{A figure drawn wider than the text block, bleeding into the margins.}
\end{figure*}

\section{Design}
\lipsum[4-6]

\appendix
\section*{Ethical Considerations}
This synthetic test paper raises no ethical concerns.

\section*{Open Science}
All artifacts of this test live in the checker repository.

\begin{thebibliography}{1}
\bibitem{esx} C.~A.\ Waldspurger. Memory resource management in VMware ESX
  server. In \emph{OSDI}, 2002.
\end{thebibliography}
\end{document}
